\section{The limiting field and its charge}\label{sec:limit}

We now prove the charge and profile assertions of
\cref{thm:priced-deficit}; \cref{prop:tf-identification} will identify
the energy coefficient.  The barrier is available at
every intermediate observation scale, before all observations are
forgotten.  This lets us select conditional fields which satisfy both
the local Thomas--Fermi equation and the small-error barrier.

\subsection{Selecting conditional data}

Take an arbitrary sequence in \cref{eq:scaling-regime}, with an
arbitrary minimizing index $n$ in each system.  By
\cref{eq:rough-deficit}, after passing to a subsequence we may suppose
\[
 q_s:=s^3(Z-n)\longrightarrow q_*.
\]
Choose integers $l\to\infty$ sufficiently slowly and a scale $j<J_*$
such that
\begin{equation}\label{eq:intermediate-scale}
 \frac1{2l}\le a_*:=\frac{r_j}{s}\le\frac1l.
\end{equation}
This is possible because $r_0/s\to0$ and the radii are dyadic.
In addition to the fixed parameters used in the barrier construction,
apply \cref{prop:inverse-comparison} at this scale with $M=l^2$, a
fixed upper height greater than $C_{\mathrm{cap}}+1$, and lower height
and tolerance sufficiently small in terms of $l$.  On its good event,
\begin{equation}\label{eq:local-tf-approximation}
 \left|\mu_j(y)-k\bigl(H'_j(y)_+\bigr)^{3/2}\right|
 \le l^{-8}|y|^{-6}
 \qquad(r_j\le|y|\le ls).
\end{equation}
Indeed, let $u=|y|^4H'_j(y)$ and
$v=(|y|^6\mu_j(y)/k)^{2/3}$.  The cap on $u$ excludes $v$ at or
above the upper height: applying the high implication at heights
approaching that endpoint would contradict the cap.  Above the
lower endpoint, the two implications at heights approaching $v$
squeeze $u$ to within the chosen tolerance of $v$.  At or below
the lower endpoint they bound $u$ above by that endpoint plus the
tolerance.  Thus $|u_+-v|\le h_l+\xi$, with both $u_+$ and $v$ in
$[0,h_h+1]$; this uses no lower bound on $u$.
Uniform continuity of $t\mapsto t^{3/2}$ on that interval proves
the assertion after choosing $h_l$ and $\xi$ sufficiently small.
The spatial range is covered because
$4Mr_j\ge2ls$.

For each fixed $l$ these are fixed parameters of
\cref{prop:inverse-comparison}; its thresholds hold eventually and
its exceptional probability tends to zero.  A slow diagonal choice
therefore makes \cref{eq:local-tf-approximation} hold outside a set
of probability tending to zero while also ensuring
\cref{eq:intermediate-scale}.

We can select data on which, in addition,
\begin{equation}\label{eq:selected-data}
 \int p_j\le1,
 \qquad
 s^3\int_{|y|>2R_0s}\mu_j(y)\,dy\le\eta_s,
 \qquad\eta_s\longrightarrow0.
\end{equation}
The first excluded probability tends to zero by
\cref{eq:barrier-error-mass}.  For the second, a master packet
centered in $|x|\le R_0s$ is contained in $|y|\le2R_0s$ for all
sufficiently far systems.  The conditional tower and
\cref{eq:exterior-cutoff} show that the expectation of its left side
tends to zero, uniformly in $j$.  Markov's inequality supplies a
deterministic threshold $\eta_s\to0$ with excluded probability
tending to zero.  These two events and the good comparison event
have positive joint probability.  Choose a point in their intersection
where all the conditional identities and barrier inequalities hold.
Only this finite-system selection is needed; every selected density
still has total mass exactly $n$.

\subsection{Compactness and the exterior boundary condition}

On the selected data, suppress the sequence index and put
\[
 F_s(x)=s^4\bigl(V-K*\mu_j\bigr)(sx),
 \qquad \rho_s(x)=s^6\mu_j(sx).
\]
Then
\[
 \Delta F_s=-4\pi Zs^3\delta_0+4\pi\rho_s,
 \qquad s^4H'_j(sx)=F_s(x)-1.
\]
The total rescaled electron mass need not be bounded.  We instead
control the screened field through its spherical mean.  Newton's
averaging identity and $\int\mu_j=n$ give
\begin{equation}\label{eq:spherical-mean}
 \overline F_s(d):=\frac1{4\pi}\int_{\mathbb S^2}F_s(d\omega)\,d\omega
 =\frac{q_s}{d}
   +\int_{|x|>d}\left(\frac1d-\frac1{|x|}\right)\rho_s(x)\,dx.
\end{equation}
In particular $\overline F_s(d)\ge q_s/d$.

On every compact annulus, \cref{prop:inverse-comparison} gives the
uniform pointwise upper bound
$F_s(x)\le1+C_{\mathrm{cap}}|x|^{-4}$ for sufficiently far systems.
Together with \cref{eq:spherical-mean} and boundedness of $q_s$, this
bounds $F_s$ in local $L^1$ away from zero: integrate
$|F_s|=2(F_s)_+-F_s$ first over spheres and then over radii.
Equation~\eqref{eq:local-tf-approximation} also bounds $\rho_s$
uniformly on these annuli.

For completeness, choose a ball whose closure avoids the origin,
and add to $F_s$ the Coulomb potential of $\rho_s$ restricted to a
slightly larger ball.  The sum is harmonic on the larger ball.
The added potential and its gradient are uniformly bounded by the
integrals of $|x-y|^{-1}$ and $|x-y|^{-2}$ against a bounded density.
Interior harmonic estimates, using the local $L^1$ bound, give
uniform bounds and equicontinuity on the smaller ball.  A diagonal
Arzel\`a--Ascoli argument now gives a subsequence converging locally
uniformly on $\R^3\setminus\{0\}$ to a field $F_\infty$.
The scaled error in \cref{eq:local-tf-approximation} tends to zero
locally uniformly, so the densities converge there to
$k((F_\infty-1)_+)^{3/2}$.  Hence
\begin{equation}\label{eq:limit-pde}
 \Delta F_\infty=4\pi k\bigl((F_\infty-1)_+\bigr)^{3/2},
 \qquad
 F_\infty(x)\le1+C_{\mathrm{cap}}|x|^{-4}.
\end{equation}
The equation holds classically off zero.  For example, bounded
local sources first give local H\"older regularity of the field;
the reaction is then locally H\"older, and interior Poisson
regularity gives the asserted conclusion.

The selected tail bound says precisely
\[
 \int_{|x|>2R_0}\rho_s(x)\,dx\longrightarrow0.
\]
Local uniform density convergence therefore makes $F_\infty$
harmonic on $|x|>2R_0$.  For $d>2R_0$, the last term of
\cref{eq:spherical-mean} is nonnegative and at most $d^{-1}$ times
this tail mass.  Consequently
\begin{equation}\label{eq:limit-exterior-mean}
 \overline F_\infty(d)=q_*/d\qquad(d>2R_0).
\end{equation}

These facts imply the uniform boundary condition
\begin{equation}\label{eq:limit-decay}
 F_\infty(x)\longrightarrow0\qquad(|x|\to\infty).
\end{equation}
Here no positivity has yet been assumed.  On large annuli rescaled
to a fixed annulus, the upper cap bounds the positive part, and
\cref{eq:limit-exterior-mean} bounds the integral of the negative
part.  Interior harmonic estimates thus give a uniform absolute
bound for $F_\infty$ throughout the exterior.  Given any sequence
$R_i\to\infty$, its dilates $F_\infty(R_i x)$ have a subsequence
converging locally on punctured space to a bounded harmonic
function, with the same absolute bound on every punctured compact
set.  Its singularity at zero is removable, and Liouville's theorem
makes it constant.  Its spherical mean is zero by
\cref{eq:limit-exterior-mean}, so the constant is zero.  If
\cref{eq:limit-decay} failed, dilating at a sequence of offending
radii and taking a convergent subsequence of the corresponding unit
directions would contradict this conclusion.  This proves uniform
decay.

\subsection{Transferring the positive barrier}

We claim that
\begin{equation}\label{eq:limit-positive-barrier}
 F_\infty(x)\ge B|x|^{-4}\qquad(x\ne0).
\end{equation}
The field in this assertion is the field without the subtracted
price.  This distinction permits the reaction comparison below.

Fix an outer radius $S>0$ and write $U_s(x)=s^4u_j(sx)$.
Inside $a_*$, the distributional inequality for $U_s$ and the
equation for $F_s$ contain the same point-nucleus and density
$\rho_s$ terms; the inequality also contains the rescaled error
from $p_j$.
Outside
$a_*$, \cref{eq:local-tf-approximation} gives, once $l>S$,
\[
 \rho_s(x)\le k\bigl((F_s(x)-1)_+\bigr)^{3/2}
                         +l^{-8}|x|^{-6}.
\]
Define the nonnegative correction potential
\begin{equation}\label{eq:barrier-correction}
 L_s=K*\left(s^6p_j(s\,\cdot)
       +l^{-8}|\cdot|^{-6}\ind_{\{a_*\le|\cdot|\le S\}}\right).
\end{equation}
Scaling \cref{eq:barrier-invariant}, subtracting the equation for
$F_s$, and using $\Delta K=-4\pi\delta_0$ gives on $B(0,S)$
\begin{equation}\label{eq:barrier-transfer-inequality}
 \Delta(U_s-F_s-L_s)
 \ge\ind_{\{|x|\ge a_*\}}\bigl[f(U_s)-f(F_s-1)\bigr].
\end{equation}
In particular, subtracting $L_s$ cancels both the accumulated
source error and the approximation error with the required sign.

Choose any
\[
 c>\left(C_{\mathrm{inv}}S^{-4}
                   -\min_{|x|=S}F_\infty(x)\right)_+.
\]
For sufficiently far systems, local uniform convergence on this
sphere, the upper bound for $U_s$, and $L_s\ge0$ give
$U_s-F_s-L_s\le c$ on the boundary.  The same inequality holds
throughout the ball.  Indeed, test
\cref{eq:barrier-transfer-inequality} with
$(U_s-F_s-L_s-c)_+$.  On its positive set,
$U_s>F_s+L_s+c>F_s-1$, so the reaction difference is nonnegative.
The test yields
\[
 \int_{B(0,S)}\left|\nabla(U_s-F_s-L_s-c)_+\right|^2\le0.
\]
It is legitimate on the full ball: the nuclear singularities in
$U_s-F_s$ cancel, the remaining offsets are continuous and locally
$H^1$, and the correction is the potential of a bounded compactly
supported density for each fixed system.  After cancellation all
right-hand sides are locally bounded, including near zero, because
the reaction term vanishes on $|x|<a_*$.  Approximation therefore
allows the compactly supported nonnegative $H^1$ test.

For each fixed $x\ne0$, $L_s(x)\to0$.  The first density in
\cref{eq:barrier-correction} has mass at most $s^3$ by
\cref{eq:selected-data} and is supported in $|x|\le a_*$.
The second has total mass at most
\[
 C l^{-8}a_*^{-3}\le C l^{-5},
\]
by \cref{eq:intermediate-scale}; its density near a fixed nonzero
evaluation point is $O_x(l^{-8})$.  Splitting its potential integral
into a small ball around that point and its complement proves the
claim.  Finally, for fixed $0<|x|<S$ and sufficiently far systems,
\[
 U_s(x)\ge B|x|^{-4}\left(1-\frac{a_*}{8|x|}\right).
\]
Pass to the limit in $U_s\le F_s+L_s+c$, then decrease $c$ to its
permitted endpoint.  Letting $S\to\infty$ and using
\cref{eq:limit-decay} proves \cref{eq:limit-positive-barrier}.

\subsection{The singular profile and uniqueness}

The dilation argument below is a direct form of the Sommerfeld comparison
principle. The homogeneous zero-price profiles are covered by
\cite[Lemma~4.4 and Remark~4.5]{Solovej2003}, including nonspherical fields.
Here dilation removes the unit price near the singularity; a separate
comparison then proves uniqueness for the priced equation.

\begin{lemma}[Unique decaying singular profile]\label{lem:singular-profile}
There is at most one classical solution $F$ of
\[
 \Delta F=4\pi k\bigl((F-1)_+\bigr)^{3/2}
       \quad\text{on }\R^3\setminus\{0\}
\]
which tends uniformly to zero at infinity and satisfies
$c|x|^{-4}\le F(x)\le C|x|^{-4}$ near zero for some $c,C>0$.
Every such solution has the uniform singular asymptotic
\begin{equation}\label{eq:sommerfeld-coefficient}
 |x|^4F(x)\longrightarrow
 A_{\mathrm S}:=\left(\frac{12}{4\pi k}\right)^2
       \qquad(|x|\downarrow0).
\end{equation}
If a solution exists, it is radial.
\end{lemma}

\begin{proof}
Let $a$ and $b$ be the lower and upper limits of $|x|^4F(x)$ at zero;
the hypotheses give $0<a\le b<\infty$.  Choose points approaching
the upper limit and put $t$ equal to their radii.  The dilates
$u_t(x)=t^4F(tx)$ satisfy
\[
 \Delta u_t=4\pi k\bigl((u_t-t^4)_+\bigr)^{3/2}.
\]
They are bounded above and below by positive multiples of $|x|^{-4}$
on compact annuli.  The local potential and harmonic compactness
argument used above supplies a subsequential classical limit $u$
with $\Delta u=4\pi k u^{3/2}$.  The definition of $b$ gives
$u(x)\le b|x|^{-4}$ for every $x\ne0$.  After taking a subsequence
of the unit directions, equality holds at a unit point.  The
difference $u-b|x|^{-4}$ has a local maximum there, so
\[
 4\pi k b^{3/2}\le12b.
\]
The analogous dilation at points approaching the lower limit touches
$a|x|^{-4}$ from above at a unit point and gives
$4\pi k a^{3/2}\ge12a$.  Since both constants are positive,
\[
 b\le\left(\frac{12}{4\pi k}\right)^2\le a.
\]
This proves \cref{eq:sommerfeld-coefficient}.  The contact argument
does not assume radiality.

Let $F_1,F_2$ be two such solutions.  For $\alpha>1$ and $\delta>0$,
the function $G=\alpha F_1+\delta$ is an upper comparison function:
\[
 \Delta G=\alpha\,4\pi k\bigl((F_1-1)_+\bigr)^{3/2}
       \le4\pi k\bigl((G-1)_+\bigr)^{3/2}.
\]
If $F_1\le1$ the inequality is immediate; otherwise use
$G-1\ge\alpha(F_1-1)$ and $\alpha^{3/2}\ge\alpha$.
The common asymptotic coefficient makes $G$ dominate $F_2$ on all
sufficiently small spheres.  Uniform decay and $\delta>0$ give
domination on all sufficiently large spheres.  On the intervening
annulus, testing with $(F_2-G)_+$ and using the nondecreasing
reaction proves $F_2\le G$.  Thus this inequality holds everywhere.
Let $\delta\downarrow0$ and $\alpha\downarrow1$, and interchange
the two solutions to obtain equality.  Rotations preserve the
equation and both boundary conditions, so uniqueness also gives
radiality.
\end{proof}

\begin{proof}[The charge assertion of \cref{thm:priced-deficit}]
The field $F_\infty$ constructed above satisfies the hypotheses of
\cref{lem:singular-profile}, by
\cref{eq:limit-pde,eq:limit-decay,eq:limit-positive-barrier}.
In particular a solution exists, is unique, and is radial.
Its exterior mean determines $q_*$ uniquely by
\cref{eq:limit-exterior-mean}.  Moreover, at any fixed $d>2R_0$,
\[
 \frac{q_*}{d}=\overline F_\infty(d)\ge Bd^{-4}>0.
\]
Denote this uniquely determined positive number by $q$.
The reaction density vanishes beyond $2R_0$, and radiality together
with \cref{eq:limit-exterior-mean} gives $F_\infty(x)=q/|x|$ there.

Every sequence of systems and every choice of minimizing indices
has bounded scaled deficits by \cref{eq:rough-deficit}.  From any
subsequence on which these deficits converge, the preceding
construction produces the same unique profile and hence the same
limit $q$.  Therefore all scaled deficits converge to $q$:
\[
 s^3(Z-n)\longrightarrow q>0
 \qquad\text{along }\cref{eq:scaling-regime}.
\]
The argument allowed an arbitrary minimizing index at every stage,
as required.  The remaining coefficient identity in
\cref{thm:priced-deficit} is proved in \cref{prop:tf-identification}.
\end{proof}
